% !TEX program = pdflatex
% Manual de uso de SUMO-GEO — compilar con: pdflatex Manual_Uso_SUMO_GEO.tex
\documentclass[11pt,a4paper]{article}

\usepackage[utf8]{inputenc}
\usepackage[T1]{fontenc}
\usepackage[margin=2.4cm]{geometry}
% (sin babel: este texlive no trae spanish.ldf; títulos en español a mano)
\AtBeginDocument{%
  \renewcommand{\contentsname}{Índice}%
  \renewcommand{\tablename}{Tabla}%
  \renewcommand{\figurename}{Figura}%
}
\usepackage{lmodern}
\usepackage{booktabs}
\usepackage{longtable}
\usepackage{array}
\usepackage{xcolor}
\usepackage{listings}
\usepackage{hyperref}
\usepackage{parskip}

\definecolor{fondo}{RGB}{245,246,248}
\definecolor{borde}{RGB}{200,205,215}
\definecolor{cmdcolor}{RGB}{30,60,110}

\lstset{
  basicstyle=\ttfamily\small,
  backgroundcolor=\color{fondo},
  frame=single, rulecolor=\color{borde},
  breaklines=true, breakatwhitespace=true,
  columns=fullflexible, keepspaces=true,
  literate={á}{{\'a}}1 {é}{{\'e}}1 {í}{{\'i}}1 {ó}{{\'o}}1 {ú}{{\'u}}1 {ñ}{{\~n}}1 {→}{{$\rightarrow$}}1 {≥}{{$\geq$}}1 {·}{{$\cdot$}}1 {°}{{$^\circ$}}1
}

\hypersetup{colorlinks=true, linkcolor=cmdcolor, urlcolor=cmdcolor}

\title{\textbf{SUMO-GEO}\\[4pt]\Large Manual de uso\\[6pt]
\large Simulación de movilidad SUMO en el navegador con MapLibre GL JS + deck.gl}
\author{Universidad de Cuenca --- Ecuador\\
\small\url{https://github.com/Pbarbecho/SUMO_GEO}}
\date{\today}

\begin{document}
\maketitle
\tableofcontents
\newpage

%=====================================================================
\section{Introducción}

\textbf{SUMO-GEO} es una aplicación web para ejecutar y visualizar simulaciones
de movilidad de \textbf{SUMO} directamente en el navegador, sin instalar ningún
software en el cliente. Ofrece edificios 3D, estimación de tráfico en tiempo
real (niveles de servicio A--F), semáforos sincronizados con las fases de SUMO,
vehículos 3D procedurales, un inspector con clic derecho y un panel de
estadísticas históricas.

\subsection{Arquitectura}

\begin{lstlisting}
SUMO (Docker) --TraCI--> Backend FastAPI --REST------> red + edificios + semaforos (GeoJSON)
 (suscripciones)                        --WebSocket--> vehiculos + LOS + fases + estadisticas
                                                        --> MapLibre GL JS + deck.gl (navegador)
\end{lstlisting}

\begin{itemize}
  \item \textbf{Simulación}: SUMO con TraCI. Cada pestaña del navegador lanza su
        propia simulación (conexiones TraCI con etiqueta única).
  \item \textbf{Backend}: FastAPI (Python). Convierte la red y los polígonos a
        GeoJSON, calcula el LOS por arista y transmite los fotogramas por WebSocket.
  \item \textbf{Frontend}: SPA estática con MapLibre GL JS (mapa base y edificios)
        y deck.gl (calles, vehículos, semáforos).
  \item \textbf{Entrega}: nginx sirve la SPA y hace de proxy inverso de
        \texttt{/api} y \texttt{/ws} (mismo origen, sin CORS).
\end{itemize}

%=====================================================================
\section{Requisitos previos}

\begin{itemize}
  \item \textbf{Docker} y \textbf{Docker Compose} (Docker Desktop en macOS/Windows,
        o \texttt{docker} + plugin \texttt{compose} en Linux).
  \item \textbf{Git} para clonar el repositorio.
  \item Conexión a internet la primera vez (descarga de imágenes Docker y teselas
        del mapa base).
  \item Opcional, solo para \emph{generar mapas nuevos} en el host:
        Python~3 con \texttt{pip3 install eclipse-sumo sumolib}.
\end{itemize}

No es necesario instalar SUMO en el equipo: el contenedor del backend lo incluye.

%=====================================================================
\section{Lanzar una simulación desde cero}

\subsection{Paso 1 --- Clonar el repositorio}

\begin{lstlisting}
git clone https://github.com/Pbarbecho/SUMO_GEO.git
cd SUMO_GEO
\end{lstlisting}

\subsection{Paso 2 --- Construir y levantar los contenedores}

\begin{lstlisting}
docker compose up --build
\end{lstlisting}

La primera ejecución tarda unos minutos (construye la imagen del backend e
instala SUMO). Las siguientes son inmediatas.

\subsection{Paso 3 --- Abrir la aplicación}

\begin{itemize}
  \item Frontend: \url{http://localhost:8082}
  \item Comprobación del backend: \url{http://localhost:8000/api/health}
\end{itemize}

Al abrir la página se conecta el WebSocket, arranca SUMO dentro del contenedor y
los vehículos comienzan a circular. Si se cambió el código o la configuración,
recargar con \emph{hard refresh}: \texttt{Cmd+Shift+R} (macOS) o
\texttt{Ctrl+Shift+R} (Windows/Linux).

\subsection{Paso 4 --- Detener}

\begin{lstlisting}
docker compose down
\end{lstlisting}

\subsection{Volver a lanzar la aplicación (segunda vez y siguientes)}

\textbf{No} hace falta \texttt{--build} en cada arranque. La imagen del backend
ya está construida y los mapas y el frontend se montan por volumen, así que
basta con:

\begin{lstlisting}
docker compose up -d
\end{lstlisting}

La regla práctica es la siguiente:

\begin{longtable}{@{}p{7.6cm}p{6.7cm}@{}}
\toprule
\textbf{Qué cambió desde la última vez} & \textbf{Comando necesario} \\
\midrule
Nada (solo volver a usar la herramienta) & \texttt{docker compose up -d} \\
Frontend (\texttt{frontend/}: HTML, JS, CSS) & Nada que reconstruir: recargar el
  navegador con \texttt{Cmd/Ctrl+Shift+R} \\
Mapas o demanda (\texttt{sumo/}, \texttt{mapa/}) & \texttt{docker compose up -d}
  y recargar el navegador (cada conexión lanza un SUMO nuevo que ya lee los
  archivos actualizados) \\
Variables \texttt{APP\_*} en \texttt{docker-compose.yml} &
  \texttt{docker compose up -d} (recrea el contenedor con el nuevo entorno) \\
Código del backend (\texttt{backend/}: Python, \texttt{requirements.txt},
  \texttt{Dockerfile}) & \texttt{docker compose up -d --build} \\
\bottomrule
\end{longtable}

En resumen: \texttt{--build} \emph{solo} cuando se toca la carpeta
\texttt{backend/}. Usarlo de más no rompe nada (Docker aprovecha la caché),
solo tarda unos segundos extra.

%=====================================================================
\section{Comandos de Docker explicados}

\subsection{\texttt{docker compose up --build}}

\begin{longtable}{@{}p{4.2cm}p{10.5cm}@{}}
\toprule
\textbf{Elemento} & \textbf{Significado} \\
\midrule
\texttt{docker compose} & Orquesta los servicios definidos en
  \texttt{docker-compose.yml} (aquí: \texttt{backend} y \texttt{frontend}). \\
\texttt{up} & Crea y arranca los contenedores, redes y volúmenes. Con la
  terminal abierta muestra los logs de ambos servicios en vivo. \\
\texttt{--build} & Fuerza la reconstrucción de la imagen del backend antes de
  arrancar. \textbf{Necesario} tras modificar \texttt{backend/} (código Python,
  \texttt{requirements.txt}, \texttt{Dockerfile}). No hace falta si solo cambian
  archivos montados por volumen (frontend, mapas, \texttt{docker-compose.yml}). \\
\texttt{-d} (opcional) & \emph{Detached}: arranca en segundo plano y devuelve la
  terminal. Ejemplo: \texttt{docker compose up -d --build}. \\
\bottomrule
\end{longtable}

\subsection{Comandos habituales}

\begin{longtable}{@{}p{6.4cm}p{8.3cm}@{}}
\toprule
\textbf{Comando} & \textbf{Uso} \\
\midrule
\texttt{docker compose up -d} & Arrancar en segundo plano (sin reconstruir). \\
\texttt{docker compose down} & Parar y eliminar los contenedores. \\
\texttt{docker compose restart backend} & Reiniciar solo el backend (por ejemplo
  tras cambiar variables de entorno no requiere \texttt{--build}, pero sí
  \texttt{up -d} para releer el YAML). \\
\texttt{docker compose logs -f} & Seguir los logs de todos los servicios
  (\texttt{Ctrl+C} para salir). \\
\texttt{docker compose logs -f backend} & Solo los logs del backend: aquí se ven
  los arranques de SUMO y los errores de TraCI. \\
\texttt{docker compose logs --tail=100 backend} & Últimas 100 líneas sin seguir. \\
\texttt{docker compose ps} & Estado de los servicios y puertos publicados. \\
\texttt{docker compose build --no-cache backend} & Reconstrucción limpia si la
  imagen quedó en mal estado. \\
\bottomrule
\end{longtable}

\subsection{Servicios y puertos}

\begin{longtable}{@{}llll@{}}
\toprule
\textbf{Servicio} & \textbf{Imagen} & \textbf{Puerto host} & \textbf{Función} \\
\midrule
\texttt{backend}  & construida de \texttt{backend/Dockerfile} & 8000 & FastAPI + SUMO \\
\texttt{frontend} & \texttt{nginx:1.27-alpine} & 8082 & SPA + proxy \texttt{/api}, \texttt{/ws} \\
\bottomrule
\end{longtable}

%=====================================================================
\section{Ubicación de los archivos de mapas de SUMO}

El repositorio incluye dos familias de escenarios. Los directorios del host se
montan \emph{de solo lectura} dentro del contenedor del backend:

\begin{longtable}{@{}lll@{}}
\toprule
\textbf{Directorio (host)} & \textbf{Montaje (contenedor)} & \textbf{Contenido} \\
\midrule
\texttt{./sumo} & \texttt{/sumo} & Centro de Cuenca (OSM) + demo sintética \\
\texttt{./mapa} & \texttt{/mapa} & Red metropolitana de Cuenca completa \\
\bottomrule
\end{longtable}

Por eso las rutas de las variables \texttt{APP\_*} empiezan por \texttt{/sumo/}
o \texttt{/mapa/}: son rutas \emph{dentro del contenedor}.

\subsection{\texttt{mapa/} --- Red metropolitana (escenario activo por defecto)}

\begin{longtable}{@{}p{5.6cm}p{8.9cm}@{}}
\toprule
\textbf{Archivo} & \textbf{Descripción} \\
\midrule
\texttt{network.net.xml} & Red vial de toda Cuenca ($\sim$93\,MB, 26$\times$19\,km,
  337 intersecciones semaforizadas). \textbf{Excluida de git} por tamaño; se
  regenera con \texttt{scripts/build\_city.sh}. \\
\texttt{metro\_low.sumocfg} / \texttt{\_mid} / \texttt{\_high} & Configuración
  SUMO de cada nivel de tráfico (Bajo / Medio / Alto). \\
\texttt{metro\_low.rou.xml} / \texttt{\_mid} / \texttt{\_high} & Demanda de 24\,h
  por nivel (6\,000 / 15\,000 / 30\,000 viajes) con perfil horario realista y
  orígenes/destinos en la componente fuertemente conexa (sin avisos
  \emph{no route}). \\
\texttt{vtypes.add.xml} & Flota de turismos (24 modelos reales: Chevrolet,
  Toyota, Kia, \ldots) con dimensiones y masas. \\
\texttt{transit.add.xml} & Tipo de vehículo \texttt{bus}. \\
\bottomrule
\end{longtable}

\subsection{\texttt{sumo/} --- Centro de Cuenca y demo}

\begin{longtable}{@{}p{5.6cm}p{8.9cm}@{}}
\toprule
\textbf{Archivo} & \textbf{Descripción} \\
\midrule
\texttt{cuenca.net.xml} & Red del centro ($\sim$1$\times$1\,km) importada de OSM,
  proyectada (se georreferencia exactamente sobre el mapa). \\
\texttt{cuenca.poly.xml} & Huellas de edificios con altura real (etiquetas OSM)
  para la extrusión 3D. \\
\texttt{cuenca.sumocfg} & Escenario demo corto (0--600\,s). \\
\texttt{cuenca\_\{low,mid,high\}.*} & Niveles de 24\,h para el centro. \\
\texttt{cuenca\_bbox.osm.xml} & Extracto OSM original de la zona. \\
\texttt{grid.net.xml}, \texttt{buildings.poly.xml}, \texttt{demo.sumocfg} &
  Demo sintética (malla 6$\times$6) generada sin internet por
  \texttt{scripts/generate\_scenario.sh}. \\
\texttt{vtypes.add.xml}, \texttt{transit.add.xml} & Copia de la flota para esta
  familia. \\
\bottomrule
\end{longtable}

\textbf{Para añadir un mapa propio}: colocar \texttt{.net.xml},
\texttt{.poly.xml} (opcional) y \texttt{.sumocfg} en \texttt{sumo/} (o
\texttt{mapa/}) y apuntar \texttt{APP\_SUMO\_CONFIG} al \texttt{.sumocfg}. Los
semáforos, la red y los edificios se derivan automáticamente al arrancar el
backend; la red debe construirse con \texttt{netconvert --tls.guess-signals}
para que existan semáforos.

%=====================================================================
\section{Configuración: variables de entorno del backend}

Se definen en la sección \texttt{environment} del servicio \texttt{backend} en
\texttt{docker-compose.yml}. Tras cambiar una variable:
\texttt{docker compose up -d} (releer YAML) y refrescar el navegador.

\begin{longtable}{@{}p{5.3cm}p{2.9cm}p{6.2cm}@{}}
\toprule
\textbf{Variable} & \textbf{Valor por defecto} & \textbf{Significado} \\
\midrule
\texttt{APP\_SUMO\_MODE} & \texttt{managed} & \texttt{managed}: el backend lanza
  SUMO. \texttt{remote}: se conecta a un servidor TraCI externo. \\
\texttt{APP\_SUMO\_CONFIG} & \texttt{/mapa/metro\_mid.sumocfg} & Escenario base:
  de él se leen red, edificios y semáforos al arrancar. \\
\texttt{APP\_SUMO\_CONFIG\_TEMPLATE} & \texttt{/mapa/metro\_\{level\}.sumocfg} &
  Plantilla del selector Bajo/Medio/Alto; \texttt{\{level\}} se sustituye por
  \texttt{low}/\texttt{mid}/\texttt{high}. \\
\texttt{APP\_SIM\_BEGIN} & \texttt{28800} & Segundo simulado de arranque
  (28\,800 = 08:00, hora punta). \\
\texttt{APP\_VIEW\_LON} / \texttt{APP\_VIEW\_LAT} & \texttt{-79.004} /
  \texttt{-2.902} & Centro inicial de la vista (útil cuando la demanda se
  concentra en una zona). \\
\texttt{APP\_NET\_FILE} / \texttt{APP\_POLY\_FILE} & --- & Sustituyen la red o
  los edificios del \texttt{.sumocfg} (p.\,ej.\ añadir \texttt{poly} a un
  escenario que no lo lista). \\
\texttt{APP\_ORIGIN\_LON} / \texttt{APP\_ORIGIN\_LAT} & Cuenca & Ancla ENU para
  redes sintéticas sin proyección (la demo \emph{grid}). \\
\texttt{APP\_STEP\_LENGTH} & \texttt{1.0} & Segundos por paso de simulación. \\
\texttt{APP\_MAX\_FPS} & \texttt{10} & Tope de fotogramas por segundo del
  WebSocket. \\
\texttt{APP\_SUMO\_HOST} / \texttt{APP\_SUMO\_PORT} & \texttt{sumo} /
  \texttt{8813} & Servidor TraCI en modo \texttt{remote}. \\
\texttt{APP\_USE\_LIBSUMO} & \texttt{false} & Usar libsumo en proceso si está
  disponible. \\
\texttt{APP\_CORS\_ORIGINS} & \texttt{*} & Orígenes CORS permitidos. \\
\bottomrule
\end{longtable}

\subsection{Cambiar de escenario}

En \texttt{docker-compose.yml} hay un bloque marcado con el escenario
\emph{activo} (red metropolitana) y, comentada, la alternativa del centro con
edificios 3D. Para cambiar: comentar unas líneas, descomentar las otras y
ejecutar \texttt{docker compose up -d}.

%=====================================================================
\section{Uso de la interfaz}

\subsection{Panel izquierdo (controles; se pliega solo, pasar el cursor para abrir)}
\begin{itemize}
  \item \textbf{Pausar / Reiniciar} la simulación.
  \item \textbf{Escenario}: nivel de tráfico Bajo / Medio / Alto (reconecta).
  \item \textbf{Velocidad}: 1--30 fps.
  \item \textbf{Edificios}, \textbf{Congestión (LOS)}, \textbf{Semáforos},
        \textbf{PoI}: capas visibles.
  \item \textbf{Calles concurridas} + umbral $\geq N$ veh: pines flotantes con el
        conteo en vivo.
\end{itemize}

\subsection{Panel derecho: Históricos}
Series temporales con eje de hora simulada: vehículos (autos/buses), CO$_2$ de la
flota, tiempo de viaje medio, espera en semáforos y desglose por tipo.

\subsection{Barra superior y cámara}
Zoom, órbita automática, alternador 2D/3D y presets de luz (amanecer / día /
atardecer / noche). El \emph{pad} \textbf{Pan-Tilt} (esquina inferior izquierda)
controla giro (horizontal) e inclinación (vertical); doble clic = norte + 3D.

\subsection{Inspector con clic derecho}
Clic derecho sobre un \textbf{vehículo} (tipo, velocidad, CO$_2$, espera,
retraso, distancia, carril, progreso de ruta), una \textbf{calle} (carriles,
límite, densidad, LOS, ocupación) o un \textbf{semáforo} (estado, tipo de
soporte). Los turismos de la flota son eléctricos (\texttt{evehicle}): su
CO$_2$ es 0; los buses sí emiten.

%=====================================================================
\section{Generación de mapas y demanda (avanzado)}

En el host, con \texttt{pip3 install eclipse-sumo sumolib}:

\begin{lstlisting}
# Importar una ciudad real de OSM (bbox = W,S,E,N):
./scripts/build_city.sh "-79.010,-2.903,-79.000,-2.893" cuenca

# Generar demanda de 24 h en 3 niveles (con duarouter):
python3 scripts/gen_traffic.py \
  --net sumo/cuenca.net.xml --types sumo/vtypes.add.xml \
  --transit sumo/transit.add.xml --outdir sumo --prefix cuenca \
  --low 3000 --mid 8000 --high 18000 --bus-share 0.05

# Variante para redes grandes (sin duarouter, sin avisos "no route"):
python3 scripts/gen_traffic.py ... --no-duarouter --connected
\end{lstlisting}

Opciones útiles de \texttt{gen\_traffic.py}: \texttt{--center-bbox="W,S,E,N"}
(concentrar la demanda en una zona), \texttt{--levels low,mid} (regenerar un
subconjunto), \texttt{--seed} (reproducibilidad).

%=====================================================================
\section{Solución de problemas}

\begin{longtable}{@{}p{6.2cm}p{8.3cm}@{}}
\toprule
\textbf{Síntoma} & \textbf{Solución} \\
\midrule
La página muestra ``esperando al backend'' & El backend aún carga la red (la
  metropolitana tarda unos segundos) o falló: revisar
  \texttt{docker compose logs -f backend}. \\
Cambios de código no se reflejan & Backend: \texttt{docker compose up -d --build}.
  Frontend: basta \emph{hard refresh} (\texttt{Cmd/Ctrl+Shift+R}). \\
``desconectado'' permanente & Ver logs del backend; si SUMO no arranca,
  verificar que \texttt{APP\_SUMO\_CONFIG} apunta a un \texttt{.sumocfg}
  existente dentro de \texttt{/sumo} o \texttt{/mapa}. \\
Puerto 8082 ocupado (\emph{port is already allocated}) & Ver qué contenedor lo usa con
  \texttt{docker ps -a | grep 8082}, detenerlo (\texttt{docker stop} +
  \texttt{docker rm}) o cambiar el mapeo \texttt{"8082:80"} del servicio
  \texttt{frontend} en \texttt{docker-compose.yml}. Cuidado: el backend publica
  el 8000 --- no asignar el mismo puerto a ambos servicios. \\
Sin semáforos en un mapa propio & La red no los trae: reconstruir con
  \texttt{netconvert --tls.guess-signals}. \\
\bottomrule
\end{longtable}

\vfill
\begin{center}
\small SUMO-GEO --- código MIT. SUMO es EPL-2.0; MapLibre GL JS y deck.gl son BSD-3/MIT.
\end{center}

\end{document}
